%%%%%%%%%%  scan idx 88  |  folio 75  %%%%%%%%%%
\textbf{LEMME 4.} \emph{On a les formules (où \(p=\codim_XY\))~:}
\begin{align}
  f^!i_!(y)&=j_!\bigl(\lambda_{-1}(\check F)y\bigr). \tag{6}\\
  \lambda_{-1}(\check F)&\equiv \gamma^{p-1}(E)\lambda^p(\check E)
    \quad\text{mod.}(1-L)\footnotemark. \tag{7}\\
  f_*(1)&=1,\quad\text{\emph{d'où} }f_*f^*=1. \tag{8}\\
  f^*i_*(\alpha)&\equiv j_*\bigl(\alpha C^{p-1}(F)\bigr)
    \quad\text{mod.}\Ker f_*\quad\text{\emph{pour tout} }\alpha\in A(Y). \tag{9}
\end{align}
\footnotetext{Cette formule inutilement compliquée peut se remplacer par exemple par VII 4.1.}

[Pour simplifier l'écriture, on a fait opérer \(K(Y)\) sur \(K(Y')\) et \(A(Y)\) sur \(A(Y')\) grâce aux homomorphismes \(g^!\) et \(g^*\). {*}Dans la formule (7), il semble que \(\gamma^{p-1}(E)\) désigne \(\lambda^{p-1}(E-2)\), c'est-à-dire \(\sum_{i+j=p-1}\binom{-p}{j}\lambda^i(E)\)\,; la congruence a lieu dans \(K(Y')\).{*}]

La formule (6) résulte de l'égalité :
\begin{equation}\tag{6 bis}
  \Tor_i^{\mathcal O_X(v)}\bigl(\mathcal O_Y(v),\mathcal O_{Y'}(v')\bigr)
  =
  \wedge^i\check F(v')
  \quad (v'\in Y',\ v=g(v')).
\end{equation}

La formule (7) résulte des formules
\begin{equation}\tag{10}
  g_!(L^{-i})=0\quad\text{pour }0<i\leq p-1,\qquad g_!(1)=1,
\end{equation}
elles-mêmes conséquences de Künneth et de la cohomologie des espaces projectifs.

La formule (8) est triviale, enfin (9) résulte de \(g_*(C^{p-1}(F))=1\) qui se déduit, par exemple, des formules immédiates :
\begin{equation}\tag{10 bis}
  g_*(\xi^i)=0\quad\text{si }0\leq i<p-1,\qquad
  g_*(\xi^{p-1})=1,
\end{equation}
avec \(\xi=C^1(L)\in A^1(Y')\).

[En fait, (9) est très vraisemblablement une identité, et pas seulement une congruence.]

